\documentclass[
  xcolor={svgnames},
  hyperref={colorlinks=true,linkcolor=blue}
]{beamer}

% Пакеты
\usepackage[utf8]{inputenc}
\usepackage[T2A]{fontenc}
\usepackage[russian]{babel}
\usepackage{graphicx}
\usepackage{tikz}
\usepackage{amsmath, amssymb}
\usepackage{beamerposter}

% Метаданные
\title{Пример постера с колонками}
\author{Егор Викторов}
\institute{РУДН}
\date{\today}

% Тема оформления
\usetheme{Boadilla}
\usecolortheme{whale}

% Размер постера (A0, портретная ориентация)
\geometry{paperwidth=841mm, paperheight=1189mm}

\begin{document}

\begin{frame}
  \maketitle

  % Первая группа: три колонки равной ширины
  \begin{columns}[T, totalwidth=\textwidth]
    \begin{column}{0.33\textwidth}
      \begin{block}{Колонка 1}
        Содержимое первой колонки (33\% ширины). \\
        Можно добавить текст, формулы, изображения. \\
        Например: $E = mc^2$.
      \end{block}
    \end{column}
    \begin{column}{0.33\textwidth}
      \begin{block}{Колонка 2}
        Содержимое второй колонки (33\% ширины). \\
        Здесь тоже можно разместить любой контент. \\
        \includegraphics[width=0.8\linewidth]{example-image-a}
      \end{block}
    \end{column}
    \begin{column}{0.33\textwidth}
      \begin{block}{Колонка 3}
        Содержимое третьей колонки (33\% ширины). \\
        Ещё немного текста для примера.
      \end{block}
    \end{column}
  \end{columns}

  % Вторая группа: две колонки разной ширины
  \vspace{1cm}
  \begin{columns}[T, totalwidth=\textwidth]
    \begin{column}{0.7\textwidth}
      \begin{block}{Широкая колонка}
        Содержимое широкой колонки (70\% ширины). \\
        Здесь можно разместить более объёмный текст или большое изображение. \\
        \begin{itemize}
          \item Пункт списка 1
          \item Пункт списка 2
          \item Пункт списка 3
        \end{itemize}
      \end{block}
    \end{column}
    \begin{column}{0.3\textwidth}
      \begin{block}{Узкая колонка}
        Содержимое узкой колонки (30\% ширины). \\
        Краткий текст или компактная таблица. \\
        \begin{tabular}{|c|c|}
          \hline
          A & B \\
          \hline
          1 & 2 \\
          \hline
        \end{tabular}
      \end{block}
    \end{column}
  \end{columns}

\end{frame}

\end{document}
